\chapter{1962:Max F. Perutz and John C. Kendrew}

\label{chap:chemistry1962}

The Nobel Prize in Chemistry for the year 1962 was awarded jointly to Max F. Perutz and John C. Kendrew.

\textbf{Motivation:}

for their studies of the structures of globular proteins

\section{Max F. Perutz}

\subsection*{Early Life and Education}

Max F. Perutz was born on 1914-05-19 in Vienna, Austria.

\subsection*{Career and Major Research}

Perutz studied chemistry at the University of Vienna and came to Cambridge in 1936, where he worked at the Cavendish Laboratory and took his doctorate in 1940. He devoted his career to the X-ray analysis of haemoglobin, developing the method of isomorphous replacement, in which heavy atoms are attached to the protein to solve the phase problem. He founded and directed the Medical Research Council Unit for Molecular Biology in Cambridge, which later became the Laboratory of Molecular Biology. In 1959 he and his colleagues obtained a structure of haemoglobin at a resolution of 5.5 angstroms.

\subsection*{Nobel Prize-winning Work}

The work for which Max F. Perutz was awarded the Nobel Prize in Chemistry in 1962 was described by the Nobel Committee as: "for their studies of the structures of globular proteins".

Perutz determined the three-dimensional structure of the globular protein haemoglobin, showing how its four subunits are arranged. The method of isomorphous replacement he developed became the key technique for determining the structures of proteins by X-ray crystallography.

\subsection*{Significance and Impact}

The solution of the haemoglobin structure made it possible to relate the behavior of a protein molecule, including its oxygen uptake and release, to its three-dimensional architecture. Perutz's technique laid the foundation for the entire field of protein crystallography.

\subsection*{Later Career and Honors}

Perutz remained at the Medical Research Council Laboratory of Molecular Biology, which he directed for many years, and continued studies of haemoglobin throughout his life. He died on 2002-02-06 in Cambridge, England.

\subsection*{Legacy}

Perutz founded protein crystallography and the Medical Research Council Laboratory of Molecular Biology, one of the most influential research institutions in modern biology, and his work opened the era of molecular structure determination for proteins.

\subsection*{Modern Developments}

Perutz's use of X-ray crystallography to solve the structure of hemoglobin founded structural biology and revealed how proteins use conformational change to function cooperatively. Hemoglobin remains a paradigm for allosteric regulation, informing drug design and the treatment of sickle-cell disease. Perutz's methods evolved into the modern structural genomics and cryo-EM era.

\subsection*{Selected Publications}

\begin{itemize}

\item \emph{Structure of Haemoglobin: A Three-Dimensional Fourier Synthesis at 5.5 {\AA} Resolution}, Proceedings of the Royal Society of London, 1960.

\end{itemize}

\clearpage

\section{John C. Kendrew}

\subsection*{Early Life and Education}

John C. Kendrew was born on 1917-03-24 in Oxford, England.

\subsection*{Career and Major Research}

Kendrew studied chemistry at Trinity College, Cambridge, and during the war worked on operational research for the Royal Air Force. He returned to Cambridge, where he joined the Medical Research Council Unit under Max F. Perutz and applied X-ray crystallography to the smaller globular protein myoglobin. In 1957 he and his colleagues computed a low-resolution structure of myoglobin, and in 1959 a structure at 2 angstroms resolution, the first atomic-level structure of a protein.

\subsection*{Nobel Prize-winning Work}

The work for which John C. Kendrew was awarded the Nobel Prize in Chemistry in 1962 was described by the Nobel Committee as: "for their studies of the structures of globular proteins".

Kendrew determined the complete three-dimensional structure of the globular protein myoglobin, showing in detail how the polypeptide chain is folded around the heme group that carries oxygen. It was the first time the atomic arrangement of a protein had been established.

\subsection*{Significance and Impact}

The myoglobin structure showed for the first time what a folded protein looks like at atomic resolution, revealing the positions of every amino acid side chain. It demonstrated that the methods of X-ray crystallography could reach atomic detail for proteins and set the pattern for the structural biology that followed.

\subsection*{Later Career and Honors}

Kendrew continued at the Medical Research Council Laboratory of Molecular Biology and later served as director-general of the European Molecular Biology Laboratory in Heidelberg. He died on 1997-08-30 in Cambridge, England.

\subsection*{Legacy}

Kendrew's myoglobin structure was the first protein structure solved at atomic resolution and stands as a founding achievement of structural biology.

\subsection*{Modern Developments}

Kendrew's determination of the three-dimensional structure of myoglobin was the first protein structure ever solved, founding structural biology and the field that today determines thousands of protein structures annually. His work demonstrated that proteins fold into unique, detailed conformations and opened the door to structure-based drug design, enzymology, and protein engineering.

\subsection*{Selected Publications}

\begin{itemize}

\item \emph{A Three-Dimensional Model of the Myoglobin Molecule Obtained by X-Ray Analysis}, Nature, 1958.

\end{itemize}

\clearpage

\section*{Chapter Summary}

The 1962 prize recognized the determination of the structures of globular proteins. Max F. Perutz solved the structure of haemoglobin and developed the method of isomorphous replacement, while John C. Kendrew used it to determine the structure of myoglobin, the first protein structure obtained at atomic resolution.

